% ===================================================================
% MT-04d: Minitest Mi colegio (Wiederholung)
% Spanisch Klasse 7
% ===================================================================
% Dauer: 20 Minuten
% Punkte: 30 BE
% ===================================================================

\documentclass[11pt, a4paper]{article}

\usepackage[utf8]{inputenc}
\usepackage[T1]{fontenc}
\usepackage[ngerman]{babel}
\usepackage{helvet}
\renewcommand{\familydefault}{\sfdefault}

\usepackage[margin=2cm]{geometry}
\usepackage{enumitem}
\usepackage{tabularx}
\usepackage{booktabs}

\usepackage{xcolor}
\usepackage{tcolorbox}
\tcbuselibrary{skins}

\usepackage{fancyhdr}
\usepackage{lastpage}

% FARBEN
\definecolor{spanisch}{HTML}{c41e3a}
\definecolor{grau}{HTML}{888888}
\definecolor{hellgrau}{HTML}{f5f5f5}

\newcommand{\fachfarbe}{spanisch}
\newcommand{\thema}{Minitest: Mi colegio (Wiederholung)}
\newcommand{\fach}{Spanisch}
\newcommand{\klasse}{7}
\newcommand{\maxpunkte}{30}
\newcommand{\zeit}{20 Minuten}
\newcommand{\datum}{\today}

\pagestyle{fancy}
\fancyhf{}
\renewcommand{\headrulewidth}{0.4pt}
\renewcommand{\footrulewidth}{0.4pt}

\lhead{\fach{} -- Klasse \klasse}
\chead{\textbf{MT-04d}}
\rhead{\datum}

\lfoot{\zeit}
\cfoot{}
\rfoot{Seite \thepage\ von \pageref{LastPage}}

\newcommand{\punktebox}[1]{\hfill\fbox{\small \_\_\_/#1}}
\newcommand{\luecke}[1]{\rule{#1}{0.4pt}}

\newtcolorbox{infobox}{
    colback=hellgrau,
    colframe=\fachfarbe,
    boxrule=1pt,
    arc=0pt,
    left=8pt, right=8pt, top=8pt, bottom=8pt
}

\newenvironment{aufgabe}[1]{%
    \vspace{12pt}
    \noindent\textbf{\textcolor{\fachfarbe}{Aufgabe #1}}
    \vspace{6pt}
    \begin{list}{}{\leftmargin=0pt}
    \item[]
}{%
    \end{list}
}

\newcommand{\es}[1]{\textcolor{\fachfarbe}{\textbf{#1}}}

\begin{document}

% ===================================================================
% KOPFBEREICH
% ===================================================================

\begin{center}
    {\Large\bfseries\textcolor{\fachfarbe}{\thema}}
\end{center}

\vspace{5pt}

\noindent
\begin{tabularx}{\textwidth}{|X|X|X|}
\hline
\textbf{Name:} \luecke{4cm} &
\textbf{Klasse:} \klasse &
\textbf{Datum:} \luecke{2cm} \\
\hline
\end{tabularx}

\vspace{10pt}

\begin{infobox}
\textbf{Zeit:} \zeit{} | \textbf{Punkte:} \maxpunkte{} BE | \textbf{Hilfsmittel:} keine
\end{infobox}

\vspace{15pt}

% ===================================================================
% AUFGABE 1: Schulfächer
% ===================================================================

\begin{aufgabe}{1 -- Schulfächer zuordnen}
Ordne die spanischen Schulfächer der deutschen Bedeutung zu. \punktebox{5}
\end{aufgabe}

\begin{center}
\begin{tabular}{rcl}
\es{las matemáticas} & \luecke{1.5cm} & Englisch \\[0.6em]
\es{la música} & \luecke{1.5cm} & Mathematik \\[0.6em]
\es{la educación física} & \luecke{1.5cm} & Geschichte \\[0.6em]
\es{la historia} & \luecke{1.5cm} & Sport \\[0.6em]
\es{el inglés} & \luecke{1.5cm} & Musik \\
\end{tabular}
\end{center}

\vspace{12pt}

% ===================================================================
% AUFGABE 2: Wochentage und Uhrzeiten
% ===================================================================

\begin{aufgabe}{2 -- Wochentage und Uhrzeiten}
Schreibe die spanische Übersetzung. \punktebox{6}
\end{aufgabe}

\noindent
a) Montag: \luecke{3cm} \hfill
b) Donnerstag: \luecke{3cm} \\[0.8em]
c) 8:00 Uhr: \luecke{4cm} \hfill
d) 9:30 Uhr: \luecke{4cm} \\[0.8em]
e) 1:00 Uhr: \luecke{4cm} \hfill
f) Samstag: \luecke{3cm} \\

\vspace{12pt}

% ===================================================================
% AUFGABE 3: Verben konjugieren
% ===================================================================

\begin{aufgabe}{3 -- Verben konjugieren}
Schreibe die richtige Form des -ar Verbs. \punktebox{7}
\end{aufgabe}

\noindent
a) (hablar, yo) Yo \luecke{2.5cm} español.\\[0.7em]
b) (estudiar, tú) Tú \luecke{2.5cm} mucho.\\[0.7em]
c) (escuchar, María) María \luecke{2.5cm} música.\\[0.7em]
d) (trabajar, nosotros) Nosotros \luecke{2.5cm} juntos.\\[0.7em]
e) (cantar, ellos) Ellos \luecke{2.5cm} en el coro.\\[0.7em]
f) (bailar, tú) Tú \luecke{2.5cm} muy bien.\\[0.7em]
g) (terminar, ella) Ella \luecke{2.5cm} a las dos.\\

\vspace{12pt}

% ===================================================================
% AUFGABE 4: Me gusta und Fragen
% ===================================================================

\begin{aufgabe}{4 -- Me gusta und Fragen}
Ergänze \es{gusta} oder \es{gustan} und ordne die Antworten zu. \punktebox{6}
\end{aufgabe}

\noindent
a) Me \luecke{2cm} el español. \hfill
b) Me \luecke{2cm} las matemáticas. \\[0.8em]
c) Me \luecke{2cm} el arte y la música.\\[1em]

\noindent
\begin{tabular}{|p{6.5cm}|p{5cm}|}
\hline
\textbf{Fragen} & \textbf{Antworten} \\
\hline
d) \es{¿A qué hora tienes español?} \luecke{1cm} & A) Los viernes a las diez. \\[0.4em]
e) \es{¿Qué día es hoy?} \luecke{1cm} & B) A las nueve. \\[0.4em]
f) \es{¿Cuándo tienes música?} \luecke{1cm} & C) Hoy es martes. \\
\hline
\end{tabular}

\vspace{12pt}

% ===================================================================
% AUFGABE 5: Sätze schreiben
% ===================================================================

\begin{aufgabe}{5 -- Schreibe auf Spanisch}
Übersetze die Sätze ins Spanische. \punktebox{6}
\end{aufgabe}

\noindent
a) Heute ist Mittwoch. \\[1.2em]
\luecke{10cm} \\[0.8em]

\noindent
b) Um 10 Uhr habe ich Musik. \\[1.2em]
\luecke{10cm} \\[0.8em]

\noindent
c) Spanisch gefällt mir. \\[1.2em]
\luecke{10cm} \\

% ===================================================================
% PUNKTEÜBERSICHT
% ===================================================================

\vspace{15pt}
\hrule
\vspace{10pt}

\noindent
\begin{tabularx}{\textwidth}{|X|c|c|c|c|c||c|}
\hline
\textbf{Aufgabe} & 1 & 2 & 3 & 4 & 5 & \textbf{Gesamt} \\
\hline
\textbf{max. Punkte} & 5 & 6 & 7 & 6 & 6 & \textbf{30} \\
\hline
\textbf{erreicht} & & & & & & \\[0.8em]
\hline
\end{tabularx}

\vspace{10pt}

\begin{flushright}
\textbf{Note:} \luecke{2cm}
\end{flushright}

\end{document}
